% GENERATED FILE. Edit canonical lessons, then run npm run generate:books.
% canonical-source-hash: fnv1a64:9bbdf85f2cd09090
% canonical-lessons: JA-C456-kyuuryou, JA-C456-kasegu, JA-C456-kyuuka, JA-C456-yatou, JA-C456-shokuba

\chapter{Pay and Time Off}
\label{ch:ja-pay-and-time-off}
\hlchaptermodality{456}

\begin{chapteropening}
\textbf{By the end of this chapter:} \emph{I can talk about pay, hiring and time off.}

Complete the last lesson of chapter 456: I can talk about pay, hiring and time off.
\end{chapteropening}

\section[\texorpdfstring{\ja{きゅうりょう} (kyūryō)}{きゅうりょう (kyūryō)}]{\ja{きゅうりょう} (kyūryō) — a salary; pay}

\label{lesson:JA-C456-kyuuryou}



\begin{quote}
Before the new one: say the Japanese for being in charge, then the Japanese for to take over (duties).
\end{quote}

\subsection*{You'll want to know: \ja{きゅうりょう}}
\textbf{\ja{きゅうりょう}} — \emph{kyūryō} — \textquotedblleft{}a salary; pay\textquotedblright{}.

\begin{quote}
\textbf{\ja{きゅうりょうが} \ja{すこし} \ja{あがりました。}} — \emph{kyūryō ga sukoshi agarimashita} — \textquotedblleft{}My salary went up a little.\textquotedblright{}
\end{quote}

\begin{cousinweb}
\textbf{\ja{きゅうりょう}} (\emph{kyūryō}) is Sino-Japanese (kango): two Chinese characters read with their on'yomi, readings borrowed from Chinese. They are \emph{kyū}, \textquotedblleft{}salary\textquotedblright{} (Middle Chinese \emph{kip}) and \emph{ryō}, \textquotedblleft{}fee\textquotedblright{} (\emph{lew}) (KANJIDIC2; Middle Chinese in Baxter's transcription).
\end{cousinweb}

\begin{grammarlens}[title={What you've built}]
One more word for describing what you lived and felt.
\end{grammarlens}

\subsection*{Your turn}
\begin{itemize}
\raggedright
  \item \emph{Say it:} \emph{kyūryō}
  \item \emph{Say it:} \emph{kyūryō}, once more
  \item \emph{Say it:} \emph{kyūryō ga sukoshi agarimashita} — My salary went up a little.
  \item \emph{Recall:} say the Japanese for being in charge, then the Japanese for to take over (duties), then say \emph{kyūryō} again
\end{itemize}

\begin{tcolorbox}[breakable,colback=teal!4,colframe=teal!35!black,title={Before you move on}]
What does \ja{きゅうりょう} mean? (\textquotedblleft{}A salary; pay\textquotedblright{}.) Say it once more.
\end{tcolorbox}
\section[\texorpdfstring{\ja{かせぐ} (kasegu)}{かせぐ (kasegu)}]{\ja{かせぐ} (kasegu) — to earn (money)}

\label{lesson:JA-C456-kasegu}



\begin{quote}
Before the new one: say the Japanese for to take over (duties), then the Japanese for a salary.
\end{quote}

\subsection*{You'll want to know: \ja{かせぐ}}
\textbf{\ja{かせぐ}} — \emph{kasegu} — \textquotedblleft{}to earn (money)\textquotedblright{}.

\begin{quote}
\textbf{\ja{あには} \ja{みせで} \ja{はたらいて、おかねを} \ja{かせいでいます。}} — \emph{ani wa mise de hataraite, okane o kaseide imasu} — \textquotedblleft{}My elder brother earns money working in a shop.\textquotedblright{}
\end{quote}

Polite \textbf{\ja{かせぎます}}, \emph{kasegimasu}; a godan verb. Past (polite): \textbf{\ja{ちちが} \ja{おかねを} \ja{かせぎました。}}, \emph{chichi ga okane o kasegimashita}.

\begin{cousinweb}
\textbf{\ja{かせぐ}} (\emph{kasegu}) is a native Japanese word (wago), inherited rather than borrowed from Chinese; its kanji spelling was attached to it later (JMdict).
\end{cousinweb}

\begin{grammarlens}[title={What you've built}]
One more word for describing what you lived and felt.
\end{grammarlens}

\subsection*{Your turn}
\begin{itemize}
\raggedright
  \item \emph{Say it:} \emph{kasegu}
  \item \emph{Say it:} \emph{kasegu}, once more
  \item \emph{Say it:} \emph{ani wa mise de hataraite, okane o kaseide imasu} — My elder brother earns money working in a shop.
  \item \emph{Recall:} say the Japanese for to take over (duties), then the Japanese for a salary, then say \emph{kasegu} again
\end{itemize}

\begin{tcolorbox}[breakable,colback=teal!4,colframe=teal!35!black,title={Before you move on}]
What does \ja{かせぐ} mean? (\textquotedblleft{}To earn (money)\textquotedblright{}.) Say it once more.
\end{tcolorbox}
\section[\texorpdfstring{\ja{きゅうか} (kyūka)}{きゅうか (kyūka)}]{\ja{きゅうか} (kyūka) — time off work; leave}

\label{lesson:JA-C456-kyuuka}



\begin{quote}
Before the new one: say the Japanese for a salary, then the Japanese for to earn (money).
\end{quote}

\subsection*{You'll want to know: \ja{きゅうか}}
\textbf{\ja{きゅうか}} — \emph{kyūka} — \textquotedblleft{}time off work; leave\textquotedblright{}.

\begin{quote}
\textbf{\ja{なつに} \ja{いっしゅうかん} \ja{きゅうかが} \ja{あります。}} — \emph{natsu ni isshūkan kyūka ga arimasu} — \textquotedblleft{}I have a week off work in summer.\textquotedblright{}
\end{quote}

\begin{cousinweb}
\textbf{\ja{きゅうか}} (\emph{kyūka}) is Sino-Japanese (kango): two Chinese characters read with their on'yomi, readings borrowed from Chinese. They are \emph{kyū}, \textquotedblleft{}rest\textquotedblright{} (Middle Chinese \emph{xjuw}) and \emph{ka}, \textquotedblleft{}leave of absence\textquotedblright{} (\emph{haeH}) (KANJIDIC2; Middle Chinese in Baxter's transcription).
\end{cousinweb}

\begin{grammarlens}[title={What you've built}]
One more word for describing what you lived and felt.
\end{grammarlens}

\subsection*{Your turn}
\begin{itemize}
\raggedright
  \item \emph{Say it:} \emph{kyūka}
  \item \emph{Say it:} \emph{kyūka}, once more
  \item \emph{Say it:} \emph{natsu ni isshūkan kyūka ga arimasu} — I have a week off work in summer.
  \item \emph{Recall:} say the Japanese for a salary, then the Japanese for to earn (money), then say \emph{kyūka} again
\end{itemize}

\begin{tcolorbox}[breakable,colback=teal!4,colframe=teal!35!black,title={Before you move on}]
What does \ja{きゅうか} mean? (\textquotedblleft{}Time off work; leave\textquotedblright{}.) Say it once more.
\end{tcolorbox}
\section[\texorpdfstring{\ja{やとう} (yatou)}{やとう (yatou)}]{\ja{やとう} (yatou) — to hire; to employ}

\label{lesson:JA-C456-yatou}



\begin{quote}
Before the new one: say the Japanese for to earn (money), then the Japanese for time off work.
\end{quote}

\subsection*{You'll want to know: \ja{やとう}}
\textbf{\ja{やとう}} — \emph{yatou} — \textquotedblleft{}to hire; to employ\textquotedblright{}.

\begin{quote}
\textbf{\ja{かいしゃは} \ja{がくせいを} \ja{ふたり} \ja{やといました。}} — \emph{kaisha wa gakusei o futari yatoimashita} — \textquotedblleft{}The company hired two students.\textquotedblright{}
\end{quote}

Polite \textbf{\ja{やといます}}, \emph{yatoimasu}; a godan verb. Past (polite): \textbf{\ja{みせが} \ja{てんいんを} \ja{やといました。}}, \emph{mise ga ten'in o yatoimashita}.

\begin{cousinweb}
\textbf{\ja{やとう}} (\emph{yatou}) is a native Japanese word (wago), inherited rather than borrowed from Chinese; its kanji spelling was attached to it later (JMdict).
\end{cousinweb}

\begin{grammarlens}[title={What you've built}]
One more word for describing what you lived and felt.
\end{grammarlens}

\subsection*{Your turn}
\begin{itemize}
\raggedright
  \item \emph{Say it:} \emph{yatou}
  \item \emph{Say it:} \emph{yatou}, once more
  \item \emph{Say it:} \emph{kaisha wa gakusei o futari yatoimashita} — The company hired two students.
  \item \emph{Recall:} say the Japanese for to earn (money), then the Japanese for time off work, then say \emph{yatou} again
\end{itemize}

\begin{tcolorbox}[breakable,colback=teal!4,colframe=teal!35!black,title={Before you move on}]
What does \ja{やとう} mean? (\textquotedblleft{}To hire; to employ\textquotedblright{}.) Say it once more.
\end{tcolorbox}
\section[\texorpdfstring{\ja{しょくば} (shokuba)}{しょくば (shokuba)}]{\ja{しょくば} (shokuba) — a workplace}

\label{lesson:JA-C456-shokuba}



\begin{quote}
Before the new one: say the Japanese for time off work, then the Japanese for to hire.
\end{quote}

\subsection*{You'll want to know: \ja{しょくば}}
\textbf{\ja{しょくば}} — \emph{shokuba} — \textquotedblleft{}a workplace\textquotedblright{}.

\begin{quote}
\textbf{\ja{しょくばの} \ja{みんなは} \ja{しんせつです。}} — \emph{shokuba no minna wa shinsetsu desu} — \textquotedblleft{}Everyone at my workplace is kind.\textquotedblright{}
\end{quote}

\begin{cousinweb}
\textbf{\ja{しょくば}} (\emph{shokuba}) is a mixed word: Sino-Japanese \emph{shoku}, \textquotedblleft{}employment\textquotedblright{} (Middle Chinese \emph{tsyik}), borrowed from Chinese, joined to native \emph{ba}, \textquotedblleft{}place\textquotedblright{}, the kun'yomi of its second kanji (KANJIDIC2; Middle Chinese in Baxter's transcription).
\end{cousinweb}

\begin{grammarlens}[title={What you've built}]
One more word for describing what you lived and felt.
\end{grammarlens}

\subsection*{Your turn}
\begin{itemize}
\raggedright
  \item \emph{Say it:} \emph{shokuba}
  \item \emph{Say it:} \emph{shokuba}, once more
  \item \emph{Say it:} \emph{shokuba no minna wa shinsetsu desu} — Everyone at my workplace is kind.
  \item \emph{Recall:} say the Japanese for time off work, then the Japanese for to hire, then say \emph{shokuba} again
\end{itemize}

\begin{tcolorbox}[breakable,colback=teal!4,colframe=teal!35!black,title={Before you move on}]
What does \ja{しょくば} mean? (\textquotedblleft{}A workplace\textquotedblright{}.) Say it once more.
\end{tcolorbox}
